\documentclass[11pt]{article}
\usepackage[utf8]{inputenc}
\usepackage{amsmath,amssymb}
\usepackage{geometry}
\usepackage{booktabs}
\usepackage{longtable}
\usepackage{siunitx}
\usepackage{enumitem}
\usepackage{hyperref}
\usepackage{float}
\usepackage{fancyhdr}
\usepackage{xcolor}
\usepackage{tcolorbox}

\geometry{margin=1in}
\pagestyle{fancy}
\fancyhead[L]{\textbf{DFD Patent \#1}}
\fancyhead[R]{\textbf{Engineering Design Document}}
\fancyfoot[C]{\thepage}
\fancyfoot[R]{\small CONFIDENTIAL --- G.T. Alcock}

\newcommand{\psi}{\psi}
\newcommand{\Reff}{\mathrm{eff}}

\definecolor{dfdblue}{RGB}{0,51,102}
\definecolor{dfdgold}{RGB}{204,153,0}

\newtcolorbox{specbox}[1]{colback=blue!5!white,colframe=dfdblue,title=#1}
\newtcolorbox{warnbox}{colback=red!5!white,colframe=red!75!black,title=\textbf{Critical Design Constraint}}
\newtcolorbox{keyresult}{colback=green!5!white,colframe=green!50!black,title=\textbf{Key Performance Result}}

\title{\Huge\textbf{$\psi$-Field Drag Reduction System}\\[0.5em]
\LARGE Engineering Design Document\\[0.5em]
\Large Patent \#1: Systems and Methods for $\psi$-Field\\
Drag Reduction and Hydrodynamic Cloaking\\[1em]
\large Document Number: DFD-ENG-P01-001\\
Revision: 1.0}

\author{Gary Thomas Alcock\\
\textit{Inventor and Principal Engineer}\\
\texttt{gary@gtacompanies.com}}

\date{March 2026}

\begin{document}

\maketitle
\thispagestyle{empty}

\vspace{2em}
\begin{center}
\begin{tabular}{ll}
\toprule
\textbf{Classification} & PROPRIETARY \\
\textbf{Patent Application} & U.S. Provisional App. 63/893,633 et al. \\
\textbf{Inventor} & Gary Thomas Alcock \\
\textbf{Date} & March 19, 2026 \\
\textbf{Status} & Preliminary Design (TRL 2--3) \\
\bottomrule
\end{tabular}
\end{center}

\newpage
\tableofcontents
\newpage

%===========================================================================
\section{Executive Summary}
%===========================================================================

This document specifies the engineering design of a $\psi$-field drag
reduction and hydrodynamic cloaking system based on Density Field Dynamics
(DFD). The system generates a controllable scalar refractive field ($\psi$)
around a vehicle hull using conformal electromagnetic (EM) transducer panels,
creating a ``refractive bubble'' that modifies the effective density,
viscosity, and compressibility of the surrounding fluid.

\begin{keyresult}
\textbf{Target performance:}
\begin{itemize}[noitemsep]
\item Drag coefficient reduction: $>60\%$ at design speed
\item Acoustic signature suppression: $>20$~dB broadband
\item Applicable media: Water, air, and trans-medium transitions
\item Response time: $<100$~ms for full bubble establishment
\item Power consumption: $<5\%$ of propulsive power savings
\end{itemize}
\end{keyresult}

The design is specified for a \textbf{30~m class underwater vehicle} as
the primary reference platform, with scaling laws provided for other
vehicle classes.

%===========================================================================
\section{System Architecture}
%===========================================================================

\subsection{Architecture Overview}

The system comprises four major subsystems:

\begin{enumerate}
\item \textbf{EM Drive Assembly (EDA)}: Conformal metamaterial transducer
  panels that generate the $\psi$ field;
\item \textbf{Sensor Suite (SS)}: Boundary-layer and far-field sensors
  providing real-time flow state estimation;
\item \textbf{Field Synthesis Controller (FSC)}: Model predictive control
  (MPC) computer that computes optimal $\psi$ distributions;
\item \textbf{Power Conditioning Unit (PCU)}: High-frequency power
  electronics driving the EM panels.
\end{enumerate}

\subsection{Functional Block Diagram}

\begin{verbatim}
+-------------------+    +------------------+
| Navigation &      |--->| Field Synthesis   |
| Mission Computer  |    | Controller (FSC)  |
+-------------------+    | [MPC Algorithm]   |
                         +--------+---------+
                                  |
                    +-------------+-------------+
                    |                           |
              +-----v------+          +---------v--------+
              | Power       |          | EM Panel Phase   |
              | Conditioning|          | & Amplitude      |
              | Unit (PCU)  |          | Controller       |
              +-----+------+          +---------+--------+
                    |                           |
              +-----v--------------------------v--------+
              |     EM Drive Assembly (EDA)              |
              |  [24 Conformal Metamaterial Panels]      |
              |  Ring 1: Panels 1-6   (Forward)         |
              |  Ring 2: Panels 7-12  (Mid-Forward)     |
              |  Ring 3: Panels 13-18 (Mid-Aft)         |
              |  Ring 4: Panels 19-24 (Aft)             |
              +-----+----------------------------------+
                    |
              +-----v----------------------------------+
              |     Refractive Bubble / Corridor        |
              |     (psi-field envelope in fluid)       |
              +-----+----------------------------------+
                    |
              +-----v----------------------------------+
              |     Sensor Suite (SS)                   |
              |  - BL pressure transducers (x48)       |
              |  - Shear stress sensors (x24)          |
              |  - Hydrophone array (x16)              |
              |  - Temperature sensors (x12)           |
              |  - IMU (x2, redundant)                 |
              +--------------+-------------------------+
                             |
                    +--------v---------+
                    | State/Flow        |
                    | Estimator         |
                    | [Kalman Filter]   |
                    +--------+---------+
                             |
                    (feedback to FSC)
\end{verbatim}

%===========================================================================
\section{EM Drive Assembly (EDA)}
\label{sec:eda}
%===========================================================================

\subsection{Panel Design}

Each EM transducer panel is a conformal metamaterial array designed to
maximize the spatial overlap integral $\mathcal{G}$ between the EM energy
density and the target $\psi$-mode profile.

\begin{specbox}{Panel Specifications}
\begin{tabular}{@{}ll@{}}
\textbf{Panel type} & Conformal metamaterial resonator array \\
\textbf{Active area} & $1.0 \times 0.5$~m (per panel) \\
\textbf{Thickness} & 25~mm (including backing structure) \\
\textbf{Unit cell type} & Split-ring resonator (SRR) + complementary \\
\textbf{Unit cell size} & $15 \times 15$~mm \\
\textbf{Unit cells per panel} & $\sim 2{,}200$ \\
\textbf{Operating frequency} & $f_0 = 1$--$100$~MHz (tunable) \\
\textbf{$2\omega$ parametric drive} & $2f_0$, phase-locked \\
\textbf{Quality factor $Q$} & $>10{,}000$ (target) \\
\textbf{Stored energy per panel} & $U_0 = 10$--$100$~kJ \\
\textbf{Modulation depth} & $m = 0.01$--$0.3$ (variable) \\
\textbf{Phase resolution} & $<0.1^\circ$ per unit cell \\
\textbf{Conformal curvature} & Matches hull radius $\pm 0.5$~mm \\
\textbf{Waterproofing} & Potted in marine-grade epoxy, IP68 \\
\textbf{Operating depth} & 0--500~m (rated to 50~bar) \\
\textbf{Mass per panel} & $\sim 35$~kg \\
\end{tabular}
\end{specbox}

\subsection{Unit Cell Design}

The metamaterial unit cell is a dual split-ring resonator (SRR) coupled
to a complementary split-ring resonator (CSRR) to achieve:
\begin{itemize}[noitemsep]
\item High EM energy density concentration ($>100\times$ free-space);
\item Controllable resonance frequency via varactor diode tuning;
\item Co-phased TE+TM superposition to restore the geometry overlap
  $\mathcal{G}$ (per the Alcock overlap restoration
  protocol~\cite{alcock2025em});
\item Dual-sector $\epsilon(\psi)/\mu(\psi)$ response for enhanced coupling.
\end{itemize}

\textbf{Materials}:
\begin{itemize}[noitemsep]
\item Copper traces (35~$\mu$m) on Rogers RT/duroid 5880 substrate
  ($\epsilon_r = 2.2$, loss tangent $<0.001$);
\item GaAs hyperabrupt varactor diodes (Skyworks SMV1234) for tuning;
\item Ferrite loading in select cells for $\mu(\psi)$ sector response.
\end{itemize}

\subsection{Panel Array Geometry}

For the 30~m reference vehicle (hull diameter $D = 3$~m):

\begin{specbox}{Array Configuration}
\begin{tabular}{@{}ll@{}}
\textbf{Number of panels} & 24 \\
\textbf{Arrangement} & 4 rings $\times$ 6 panels per ring \\
\textbf{Ring spacing} & 5~m (along hull axis) \\
\textbf{Ring 1 position} & 2.5~m aft of nose \\
\textbf{Ring 4 position} & 17.5~m aft of nose \\
\textbf{Coverage angle} & $6 \times 60^\circ = 360^\circ$ per ring \\
\textbf{Total active area} & 12~m$^2$ \\
\textbf{Total panel mass} & 840~kg \\
\textbf{Target bubble radius} & 1.5$D$ = 4.5~m \\
\textbf{Forward corridor length} & 2.0$D$ = 6.0~m \\
\end{tabular}
\end{specbox}

\subsection{Phase Control Strategy}

The 24 panels are phase-locked to synthesize a target $\psi$ envelope
with the following features:
\begin{itemize}[noitemsep]
\item \textbf{Axial symmetry}: Panels in each ring driven with equal
  amplitude and progressive phase to create a cylindrical $\psi$ shell;
\item \textbf{Axial gradient}: Forward rings driven at higher amplitude
  to project a leading corridor;
\item \textbf{$2\omega$ parametric component}: All panels carry a
  phase-locked $2\omega$ modulation to drive the parametric channel;
\item \textbf{Adaptive beam steering}: Phase offsets between panels allow
  the bubble center to track vehicle motion and compensate for crossflow.
\end{itemize}

The phase pattern for ring $j$, panel $k$ at time $t$ is:
\begin{equation}
\phi_{jk}(t) = \omega t + \Delta\phi_j^\mathrm{axial}
+ \frac{2\pi k}{6}\,n_\mathrm{azimuth}
+ \phi_{jk}^\mathrm{MPC}(t),
\end{equation}
where $\Delta\phi_j^\mathrm{axial}$ controls the corridor projection,
$n_\mathrm{azimuth}$ the azimuthal mode order (typically 0 for symmetric
operation), and $\phi_{jk}^\mathrm{MPC}$ the real-time MPC correction.

%===========================================================================
\section{Sensor Suite}
\label{sec:sensors}
%===========================================================================

\begin{specbox}{Sensor Inventory}
\begin{longtable}{@{}llll@{}}
\toprule
\textbf{Sensor Type} & \textbf{Qty} & \textbf{Specification} & \textbf{Location} \\
\midrule
\endhead
Pressure transducer & 48 & Kulite XTL-190, 0--50~bar, & Hull surface, \\
 & & 100~kHz bandwidth & 12 per ring \\
\addlinespace
Shear stress sensor & 24 & MEMS floating element, & Hull surface, \\
 & & 0.01~Pa resolution & 6 per ring \\
\addlinespace
Hydrophone & 16 & B\&K 8103, 0.1~Hz--180~kHz, & Flush-mounted, \\
 & & $-211$~dB re 1V/$\mu$Pa & 4 per ring \\
\addlinespace
Temperature sensor & 12 & Pt100 RTD, 0.01~K resolution & Hull surface \\
\addlinespace
IMU & 2 & Honeywell HG9900, & Internal, \\
 & & 0.003$^\circ$/hr drift & midship \\
\addlinespace
Speed log & 1 & Doppler, 0.01~m/s resolution & Keel \\
\addlinespace
Depth sensor & 2 & Paroscientific Digiquartz, & Bow, stern \\
 & & 0.01\% accuracy & \\
\addlinespace
Acoustic Doppler & 1 & 300~kHz ADCP, & Bow \\
profiler & & 128-beam & \\
\bottomrule
\end{longtable}
\end{specbox}

\subsection{Sensor Placement Strategy}

Sensors are concentrated at four diagnostic stations:
\begin{enumerate}[noitemsep]
\item \textbf{Nose station} ($x = 1$~m): Stagnation pressure, freestream
  conditions;
\item \textbf{Forward bubble boundary} ($x = 5$~m): Transition detection,
  bubble onset confirmation;
\item \textbf{Mid-hull station} ($x = 15$~m): Fully developed bubble
  diagnostics;
\item \textbf{Aft closure station} ($x = 25$~m): Wake and cavity closure
  monitoring.
\end{enumerate}

\subsection{Data Acquisition}

\begin{specbox}{DAQ Specifications}
\begin{tabular}{@{}ll@{}}
\textbf{Sampling rate} & 500~kHz per channel (pressure, hydrophones) \\
\textbf{ADC resolution} & 24-bit \\
\textbf{Total channels} & 105 \\
\textbf{Data rate} & $\sim 150$~MB/s \\
\textbf{Interface} & EtherCAT, 1~$\mu$s synchronization \\
\textbf{Storage} & 4~TB NVMe RAID \\
\textbf{Processing} & FPGA front-end (Xilinx Kintex UltraScale+) \\
\end{tabular}
\end{specbox}

%===========================================================================
\section{Field Synthesis Controller (FSC)}
\label{sec:fsc}
%===========================================================================

\subsection{Control Architecture}

The FSC implements a three-layer control hierarchy:

\begin{enumerate}
\item \textbf{Strategic layer} (1~Hz update): Mission-level objectives
  (speed, stealth mode, energy budget) set the target $\psi$ envelope
  parameters;
\item \textbf{Tactical layer} (100~Hz update): MPC optimizer computes
  optimal panel amplitudes and phases to achieve the target envelope
  given current flow state;
\item \textbf{Reactive layer} (10~kHz update): Fast feedback loop
  stabilizes the bubble against transient disturbances (gusts, waves,
  maneuvers).
\end{enumerate}

\subsection{MPC Formulation}

The MPC cost function to be minimized at each tactical update is:
\begin{equation}
J = \int_0^{T_h}\!\left[w_D\,F_D^2 + w_T\,\|\tau'\|^2
+ w_A\,\|p_\mathrm{ac}\|^2
+ w_P\,P_\mathrm{EM}^2\right]dt
+ w_\Delta\,\|\Delta\boldsymbol{\phi}\|^2,
\end{equation}
where:
\begin{itemize}[noitemsep]
\item $F_D$: estimated drag force;
\item $\tau'$: turbulence intensity (from shear sensors);
\item $p_\mathrm{ac}$: far-field acoustic pressure (from hydrophones);
\item $P_\mathrm{EM}$: EM power consumption;
\item $\Delta\boldsymbol{\phi}$: phase change vector (penalizes rapid
  actuation);
\item $w_D, w_T, w_A, w_P, w_\Delta$: tunable weights;
\item $T_h = 0.5$~s: prediction horizon.
\end{itemize}

\subsection{State Estimator}

An unscented Kalman filter (UKF) estimates the full flow state from sensor
data:
\begin{equation}
\hat{\mathbf{x}} = [\hat{C}_D,\;\hat{\delta}_{BL},\;
\hat{v}_\infty,\;\hat{\alpha},\;\hat{\beta},\;
\hat{\psi}_1\ldots\hat{\psi}_{N_m}]^T,
\end{equation}
where the last $N_m$ elements are the estimated $\psi$-mode amplitudes
(observable through their effect on pressure and acoustic sensors).

\subsection{Computing Hardware}

\begin{specbox}{FSC Hardware}
\begin{tabular}{@{}ll@{}}
\textbf{Processor} & NVIDIA Jetson AGX Orin (275~TOPS AI) \\
\textbf{FPGA co-processor} & Xilinx Kintex UltraScale+ KU15P \\
\textbf{MPC solver} & OSQP (real-time convex optimization) \\
\textbf{UKF implementation} & Custom CUDA kernels \\
\textbf{Latency} & $<2$~ms sensor-to-actuator \\
\textbf{Redundancy} & Dual-redundant, hot standby \\
\textbf{Power consumption} & 150~W (both units) \\
\textbf{Enclosure} & Pressure-rated, conduction-cooled \\
\end{tabular}
\end{specbox}

%===========================================================================
\section{Power Conditioning Unit (PCU)}
\label{sec:pcu}
%===========================================================================

\subsection{Power Architecture}

The PCU converts vehicle bus power (typically 440~VAC, 60~Hz or 300~VDC)
to the high-frequency, phase-controlled drive signals required by the EM
panels.

\begin{specbox}{PCU Specifications}
\begin{tabular}{@{}ll@{}}
\textbf{Input power} & 440~VAC 3-phase or 300~VDC bus \\
\textbf{Total output power} & 500~kW peak, 100~kW continuous \\
\textbf{Output channels} & 24 (one per panel) \\
\textbf{Per-channel power} & 21~kW peak, 4.2~kW continuous \\
\textbf{Frequency range} & 1--200~MHz \\
\textbf{$2\omega$ generation} & Phase-locked loop, $<0.01^\circ$ jitter \\
\textbf{Amplitude resolution} & 14-bit DAC per channel \\
\textbf{Phase resolution} & $0.02^\circ$ \\
\textbf{Efficiency} & $>92\%$ (GaN power stage) \\
\textbf{Topology} & Class-E/F GaN HEMT amplifiers \\
\textbf{Cooling} & Liquid-cooled cold plates \\
\textbf{Mass} & 450~kg \\
\textbf{Volume} & $0.8 \times 0.6 \times 0.5$~m \\
\end{tabular}
\end{specbox}

\subsection{Waveform Generation}

Each channel generates a composite waveform:
\begin{equation}
V_k(t) = V_0\,[1 + m\cos(2\omega t + \phi_k^{2\omega})]
\cdot\cos(\omega t + \phi_k),
\end{equation}
where $V_0$ is the carrier amplitude, $m$ the parametric modulation depth,
$\phi_k$ the panel phase, and $\phi_k^{2\omega}$ the parametric phase.
The carrier at $\omega$ provides the base EM energy storage; the $2\omega$
envelope drives the parametric $\psi$ channel.

\subsection{Safety Interlocks}

\begin{warnbox}
The EM panels store significant energy ($U_0 \sim 100$~kJ per panel,
2.4~MJ total). Safety interlocks include:
\begin{enumerate}[noitemsep]
\item Emergency discharge resistor bank (10~s full discharge);
\item Over-temperature shutdown ($>85^\circ$C panel surface);
\item Over-pressure shutdown (if hull sensor detects breach);
\item $\psi$-amplitude limiter (prevents exceeding $|\delta\psi|_\mathrm{max}$);
\item Automatic shutdown on loss of FSC communication ($>10$~ms).
\end{enumerate}
\end{warnbox}

%===========================================================================
\section{Bill of Materials}
\label{sec:bom}
%===========================================================================

\begin{longtable}{@{}p{0.5cm}p{5cm}p{1.5cm}p{2cm}p{2cm}p{2.5cm}@{}}
\caption{Bill of Materials --- $\psi$-Field Drag Reduction System
(30~m reference vehicle)} \\
\toprule
\textbf{\#} & \textbf{Item} & \textbf{Qty} & \textbf{Unit Cost}
& \textbf{Total Cost} & \textbf{Supplier} \\
\midrule
\endfirsthead
\toprule
\textbf{\#} & \textbf{Item} & \textbf{Qty} & \textbf{Unit Cost}
& \textbf{Total Cost} & \textbf{Supplier} \\
\midrule
\endhead
\multicolumn{6}{c}{\textit{EM Drive Assembly}} \\
\midrule
1 & SRR/CSRR metamaterial panel assembly (1.0$\times$0.5~m) & 24
& \$45{,}000 & \$1{,}080{,}000 & Custom fab \\
2 & GaAs varactor diodes (SMV1234, 2200/panel) & 52{,}800
& \$1.20 & \$63{,}360 & Skyworks \\
3 & Rogers RT/duroid 5880 substrate (18$\times$24 in sheets) & 200
& \$380 & \$76{,}000 & Rogers Corp \\
4 & Marine-grade epoxy potting compound (kg) & 500
& \$45 & \$22{,}500 & 3M Scotchcast \\
5 & Panel mounting hardware (titanium) & 24 sets
& \$2{,}800 & \$67{,}200 & Custom CNC \\
6 & RF interconnect cables (low-loss, 3~m) & 48
& \$450 & \$21{,}600 & Times Microwave \\
\midrule
\multicolumn{6}{c}{\textit{Sensor Suite}} \\
\midrule
7 & Kulite XTL-190 pressure transducer & 48
& \$1{,}200 & \$57{,}600 & Kulite \\
8 & MEMS shear stress sensor (custom) & 24
& \$3{,}500 & \$84{,}000 & Custom MEMS \\
9 & B\&K 8103 hydrophone & 16
& \$4{,}200 & \$67{,}200 & Br\"uel \& Kj\ae r \\
10 & Pt100 RTD temperature sensor & 12
& \$85 & \$1{,}020 & Omega Eng. \\
11 & Honeywell HG9900 IMU & 2
& \$95{,}000 & \$190{,}000 & Honeywell \\
12 & Doppler speed log & 1
& \$28{,}000 & \$28{,}000 & RDI Teledyne \\
13 & Paroscientific Digiquartz depth sensor & 2
& \$6{,}500 & \$13{,}000 & Paroscientific \\
14 & ADCP (300~kHz, 128-beam) & 1
& \$42{,}000 & \$42{,}000 & RDI Teledyne \\
\midrule
\multicolumn{6}{c}{\textit{Field Synthesis Controller}} \\
\midrule
15 & NVIDIA Jetson AGX Orin & 2
& \$1{,}999 & \$3{,}998 & NVIDIA \\
16 & Xilinx KU15P FPGA board & 2
& \$12{,}000 & \$24{,}000 & AMD/Xilinx \\
17 & EtherCAT master/slave modules & 8
& \$850 & \$6{,}800 & Beckhoff \\
18 & NVMe SSD (4~TB) & 2
& \$450 & \$900 & Samsung \\
19 & Pressure-rated enclosure & 2
& \$8{,}500 & \$17{,}000 & Custom \\
\midrule
\multicolumn{6}{c}{\textit{Power Conditioning Unit}} \\
\midrule
20 & GaN HEMT Class-E/F amplifier module & 24
& \$8{,}500 & \$204{,}000 & Wolfspeed/Custom \\
21 & Phase-locked loop synthesizer & 1
& \$15{,}000 & \$15{,}000 & Analog Devices \\
22 & DC-DC converter (bus to amplifier) & 24
& \$1{,}200 & \$28{,}800 & Vicor \\
23 & Liquid cooling system (cold plates + pump) & 1
& \$18{,}000 & \$18{,}000 & Lytron \\
24 & Emergency discharge resistor bank & 1
& \$5{,}000 & \$5{,}000 & Ohmite \\
25 & PCU enclosure (pressure-rated) & 1
& \$12{,}000 & \$12{,}000 & Custom \\
\midrule
\multicolumn{6}{c}{\textit{Integration and Test}} \\
\midrule
26 & Cabling harness (vehicle-specific) & 1 lot
& --- & \$85{,}000 & Custom \\
27 & Integration labor (est. 2000 hrs) & ---
& --- & \$400{,}000 & --- \\
28 & Water tunnel test campaign (200 hrs) & ---
& --- & \$250{,}000 & --- \\
29 & Software development \& V\&V & ---
& --- & \$350{,}000 & --- \\
\midrule
\midrule
& \textbf{TOTAL ESTIMATED COST} & & & \textbf{\$3{,}283{,}978} & \\
\bottomrule
\end{longtable}

%===========================================================================
\section{Performance Metrics and Predictions}
\label{sec:performance}
%===========================================================================

\subsection{Drag Reduction Performance}

Performance predictions are parameterized by the EM--$\psi$ coupling
strength $|\lambda - 1|$, which remains the critical unknown.

\begin{specbox}{Performance vs.\ Coupling Strength}
\begin{tabular}{@{}lcccc@{}}
\toprule
$|\lambda - 1|$ & $\delta\psi$ & $\mathcal{R}_D$ (\%)
& Acoustic (dB) & Power (kW) \\
\midrule
$10^{-5}$ (accidental bound) & $10^{-12}$ & $<0.001$ & $<0.01$
& 0.5 \\
$10^{-4}$ & $10^{-9}$ & $<0.01$ & 0.1 & 2 \\
$10^{-3}$ & $10^{-6}$ & 0.3 & 3 & 15 \\
$10^{-2}$ & $10^{-3}$ & 30 & 15 & 50 \\
$10^{-1}$ & $0.1$ & 74 & 25 & 100 \\
$\sim 1$ & $0.3$ & $>90$ & $>30$ & 200 \\
\bottomrule
\end{tabular}
\end{specbox}

\begin{warnbox}
For $|\lambda - 1| \leq 10^{-5}$ (current accidental bound), the drag
reduction is negligible with this system design. The system becomes
practically useful only if $|\lambda - 1| \geq 10^{-3}$, which has
\emph{not yet been excluded} by any measurement. The first priority is
therefore the $\lambda$ measurement campaign (Section~\ref{sec:test_campaign}).
\end{warnbox}

\subsection{Speed Envelope}

For a 30~m underwater vehicle at $|\lambda - 1| \sim 0.1$:

\begin{specbox}{Speed Envelope}
\begin{tabular}{@{}lcccc@{}}
\toprule
\textbf{Parameter} & \textbf{Baseline} & \textbf{$\psi$-Active}
& \textbf{Improvement} \\
\midrule
Maximum speed (kts) & 25 & 45+ & $80\%$ \\
Cruise speed (kts) & 12 & 12 & --- \\
Cruise power (kW) & 500 & 180 & $64\%$ savings \\
Range at cruise & 5000~nm & 13{,}900~nm & $178\%$ \\
Acoustic signature & 120~dB re 1$\mu$Pa & $<95$~dB & $25$~dB \\
\bottomrule
\end{tabular}
\end{specbox}

\subsection{Bubble Dynamics}

\begin{specbox}{Bubble Characteristics}
\begin{tabular}{@{}ll@{}}
\textbf{Bubble radius} & 1.5$D$ = 4.5~m (radial) \\
\textbf{Forward corridor length} & 2.0$D$ = 6.0~m \\
\textbf{Aft wake suppression} & 1.0$D$ = 3.0~m \\
\textbf{Establishment time} & $<100$~ms ($\sim 5\,\Omega_\psi^{-1}$) \\
\textbf{Tracking bandwidth} & 10~Hz (maneuver following) \\
\textbf{Crossflow tolerance} & $\pm 15^\circ$ angle of attack \\
\textbf{Depth range} & 0--500~m (pressure-compensated) \\
\textbf{$\psi$ uniformity (radial)} & $\pm 5\%$ within bubble \\
\textbf{Gradient smoothness} & $|\nabla^2\psi|/|\nabla\psi|^2 < 0.1$ \\
\end{tabular}
\end{specbox}

%===========================================================================
\section{Operating Modes}
\label{sec:modes}
%===========================================================================

\subsection{Mode 1: Maximum Drag Reduction}

All 24 panels at full power, symmetric bubble, forward corridor extended.
Optimized for maximum speed or minimum power at cruise.
\begin{itemize}[noitemsep]
\item Panels: 100\% amplitude, symmetric phase
\item $2\omega$ drive: Maximum modulation depth ($m = 0.3$)
\item MPC objective: $w_D = 1$, $w_A = 0$, $w_P = 0.01$
\item Power: 100--200~kW continuous
\end{itemize}

\subsection{Mode 2: Acoustic Stealth}

Bubble shaped for maximum acoustic cloaking, drag reduction secondary.
\begin{itemize}[noitemsep]
\item Panels: Asymmetric amplitude (stronger radial gradient)
\item $2\omega$ drive: Moderate ($m = 0.1$)
\item MPC objective: $w_D = 0.1$, $w_A = 1$, $w_P = 0.1$
\item Power: 50--100~kW continuous
\item Achieves: $>25$~dB broadband suppression over $2\pi$ steradians
\end{itemize}

\subsection{Mode 3: Trans-Medium Transition}

Rapid reconfiguration for water-to-air or air-to-water passage.
\begin{itemize}[noitemsep]
\item Panels: Maximum power, corridor compressed to $0.5D$ forward
\item $2\omega$ drive: Maximum ($m = 0.3$), frequency swept during transition
\item MPC objective: Minimize transition force impulse
\item Power: 500~kW peak (10~s burst)
\item Sequence: Pre-condition bubble $\to$ transition $\to$ air-mode reconfigure
\end{itemize}

\subsection{Mode 4: Idle / Measurement}

Low-power mode for $\lambda$ characterization and sensor calibration.
\begin{itemize}[noitemsep]
\item Panels: 10\% amplitude, sequential activation
\item $2\omega$ drive: Swept frequency
\item Purpose: Measure $\psi$-mode response, estimate $|\lambda - 1|$
\item Power: 5~kW
\end{itemize}

%===========================================================================
\section{Power Budget}
\label{sec:power}
%===========================================================================

\begin{specbox}{System Power Budget (Mode 1 --- Maximum Drag Reduction)}
\begin{tabular}{@{}lrr@{}}
\toprule
\textbf{Subsystem} & \textbf{Peak (kW)} & \textbf{Continuous (kW)} \\
\midrule
EM panels (24 $\times$ carrier) & 250 & 50 \\
EM panels ($2\omega$ modulation) & 250 & 50 \\
PCU losses (8\% inefficiency) & 40 & 8 \\
FSC computing & 0.15 & 0.15 \\
Sensor suite \& DAQ & 0.3 & 0.3 \\
Cooling system & 15 & 8 \\
\midrule
\textbf{Total} & \textbf{555} & \textbf{116} \\
\midrule
Propulsive power saved (64\%) & --- & $-320$ \\
\textbf{Net power change} & --- & \textbf{$-204$} \\
\bottomrule
\end{tabular}
\end{specbox}

\begin{keyresult}
At the design coupling strength ($|\lambda - 1| \sim 0.1$), the system
saves 204~kW net---the drag reduction power savings (320~kW) exceed the
system's own consumption (116~kW) by a factor of 2.8. The system is
\textbf{net energy positive} under these conditions.
\end{keyresult}

%===========================================================================
\section{Test Campaign Plan}
\label{sec:test_campaign}
%===========================================================================

\subsection{Phase 0: $\lambda$ Measurement (Critical Path)}

Before any drag reduction testing, the EM--$\psi$ coupling parameter
$|\lambda - 1|$ must be measured or further constrained.

\begin{specbox}{Phase 0 --- $\lambda$ Measurement}
\begin{tabular}{@{}ll@{}}
\textbf{Duration} & 6 months \\
\textbf{Location} & Laboratory (no vehicle required) \\
\textbf{Equipment} & High-$Q$ microwave cavity, PLL, interferometer \\
\textbf{Method} & Parametric $2\omega$ drive per Alcock (2025) protocol \\
\textbf{Target sensitivity} & $|\lambda - 1| \leq 10^{-10}$ \\
\textbf{Decision gate} & If $|\lambda - 1| > 10^{-4}$: proceed to Phase 1 \\
 & If $|\lambda - 1| < 10^{-10}$: project reassessment \\
\textbf{Estimated cost} & \$150{,}000 \\
\end{tabular}
\end{specbox}

\subsection{Phase 1: Scale Model Water Tunnel Test}

\begin{specbox}{Phase 1 --- Scale Model}
\begin{tabular}{@{}ll@{}}
\textbf{Duration} & 12 months \\
\textbf{Scale} & 1:10 (3~m model) \\
\textbf{Panels} & 8 (2 rings $\times$ 4) \\
\textbf{Test facility} & Large cavitation tunnel \\
\textbf{Test matrix} & Speed $\times$ EM power $\times$ frequency \\
\textbf{Key measurements} & $C_D$, BL profiles, acoustic field \\
\textbf{Success criterion} & Measurable $\Delta C_D$ correlating with \\
 & $2\omega$ drive (not DC MHD) \\
\textbf{Estimated cost} & \$1{,}200{,}000 \\
\end{tabular}
\end{specbox}

\subsection{Phase 2: Full-Scale Prototype}

\begin{specbox}{Phase 2 --- Full-Scale}
\begin{tabular}{@{}ll@{}}
\textbf{Duration} & 24 months \\
\textbf{Platform} & 30~m unmanned underwater vehicle \\
\textbf{Full system} & Per this design document \\
\textbf{Sea trials} & 6-month campaign \\
\textbf{Key measurements} & Range extension, speed increase, \\
 & acoustic signature reduction \\
\textbf{Estimated cost} & \$8{,}000{,}000 (excluding vehicle) \\
\end{tabular}
\end{specbox}

\subsection{Phase 3: Multi-Platform Demonstration}

\begin{specbox}{Phase 3 --- Demonstration}
\begin{tabular}{@{}ll@{}}
\textbf{Duration} & 18 months \\
\textbf{Platforms} & Surface vessel, aircraft, trans-medium body \\
\textbf{Objective} & Demonstrate universality across media \\
\textbf{Estimated cost} & \$15{,}000{,}000 \\
\end{tabular}
\end{specbox}

%===========================================================================
\section{Risk Assessment}
\label{sec:risk}
%===========================================================================

\begin{longtable}{@{}p{4cm}p{1.5cm}p{1.5cm}p{6cm}@{}}
\caption{Risk Register} \\
\toprule
\textbf{Risk} & \textbf{Likelihood} & \textbf{Impact} & \textbf{Mitigation} \\
\midrule
\endfirsthead
\midrule
\endhead
$|\lambda - 1| = 0$ exactly & Low & Critical &
Phase 0 measurement; if confirmed, project terminates cleanly.
DFD theory allows $|\lambda - 1| \neq 0$ but does not require it. \\
\addlinespace
$|\lambda - 1| < 10^{-4}$ & Medium & High &
Insufficient for practical drag reduction with current panel technology.
Mitigation: develop higher-$Q$ metamaterials, larger arrays, resonant
enhancement. \\
\addlinespace
Panel $Q$ degradation in seawater & Medium & Medium &
Marine fouling, corrosion reduce $Q$. Mitigation: anti-fouling coatings,
hermetic potting, periodic $Q$ monitoring with automatic frequency retune. \\
\addlinespace
Bubble instability at high speed & Medium & Medium &
Shear-driven bubble distortion above design speed. Mitigation: MPC
adaptive control with 10~Hz tracking bandwidth. \\
\addlinespace
EMI with vehicle systems & Low & Medium &
High-power RF near sensitive electronics. Mitigation: shielded cable runs,
frequency coordination, EMC testing. \\
\addlinespace
Regulatory / classification & Low & Medium &
Novel technology may face classification society resistance. Mitigation:
early engagement with DNV, Lloyd's, ABS. \\
\bottomrule
\end{longtable}

%===========================================================================
\section{Scaling Laws for Other Platforms}
\label{sec:scaling}
%===========================================================================

The system scales to other vehicle classes as follows:

\begin{specbox}{Scaling Relations}
\begin{tabular}{@{}ll@{}}
\textbf{Panel count} & $N_p \propto (L/L_\mathrm{ref})^2$ \\
\textbf{Bubble radius} & $R_b \propto D$ \\
\textbf{Power} & $P \propto R_b^3 \propto D^3$ \\
\textbf{Panel area} & $A_\mathrm{panel} \propto D^2$ \\
\textbf{Response time} & $\tau \propto R_b / c_s \propto D$ \\
\end{tabular}
\end{specbox}

\begin{specbox}{Platform Scaling Examples}
\begin{tabular}{@{}lccccc@{}}
\toprule
\textbf{Platform} & $L$ (m) & $D$ (m) & $N_p$ & $P$ (kW)
& Cost (est.) \\
\midrule
Small UUV & 3 & 0.5 & 4 & 2 & \$200K \\
Large UUV & 10 & 1.5 & 8 & 15 & \$800K \\
\textbf{Reference sub} & \textbf{30} & \textbf{3} & \textbf{24}
& \textbf{116} & \textbf{\$3.3M} \\
Surface vessel (50~m) & 50 & 8 & 48 & 500 & \$8M \\
Aircraft (30~m span) & 30 & 4 & 24 & 80 & \$4M \\
Trans-medium (5~m) & 5 & 1 & 6 & 5 & \$400K \\
\bottomrule
\end{tabular}
\end{specbox}

%===========================================================================
\section{Intellectual Property Coverage}
\label{sec:ip}
%===========================================================================

This design is covered by the following patent applications:
\begin{enumerate}[noitemsep]
\item U.S. Provisional App. 63/865,279 --- DFD Gravitational and Signal
  Correction
\item U.S. Provisional App. 63/865,283 --- DFD Energy Transmission
\item U.S. Provisional App. 63/865,293 --- DFD Information Processing
\item U.S. Provisional App. 63/868,073 --- Artificial Generation of Scalar
  Density Fields
\item U.S. Provisional App. 63/893,633 --- EM Coupling, Actuation, and
  Application of SRF
\item Non-provisional: ``Systems and Methods for $\psi$-Field Drag Reduction
  and Hydrodynamic Cloaking'' (October 2025)
\end{enumerate}

Key claims covered: apparatus (Claim 1), method (Claim 2), cloaking system
(Claim 3), MPC control, metamaterial panels, parametric drive, trans-medium
operation, acoustic suppression.

%===========================================================================
\section{Conclusion and Recommendations}
\label{sec:conclusion}
%===========================================================================

The $\psi$-field drag reduction system is a technically feasible design
contingent on one critical parameter: the EM--$\psi$ coupling strength
$|\lambda - 1|$. The recommended path forward is:

\begin{enumerate}
\item \textbf{Immediate} (0--6 months): Execute Phase 0 $\lambda$ measurement.
  This requires only \$150K and existing laboratory equipment. The result
  determines the entire program's viability.

\item \textbf{Near-term} (6--18 months): If $|\lambda - 1| > 10^{-4}$,
  proceed to Phase 1 scale-model testing. Design and fabricate 8-panel
  array for water tunnel installation.

\item \textbf{Medium-term} (18--42 months): Phase 2 full-scale prototype
  on a 30~m UUV platform. Target: first-ever demonstration of
  $\psi$-field drag reduction in open water.

\item \textbf{Long-term} (42--60 months): Phase 3 multi-platform
  demonstration including trans-medium capability. Technology transition
  to maritime, defense, and aerospace applications.
\end{enumerate}

The potential impact is transformative: a net-energy-positive drag reduction
system applicable to \emph{any} fluid medium, with simultaneous acoustic
cloaking capability, would fundamentally change the economics and
capabilities of underwater, surface, and aerospace vehicles.

\begin{thebibliography}{10}
\bibitem{alcock2025dfd}
G.~Alcock, \emph{Density Field Dynamics: A Complete Unified Theory},
v3.2, March 2026.

\bibitem{alcock2025em}
G.~Alcock, \emph{Accidental and Intentional Constraints on an
EM$\to\psi$ Back-Reaction Coupling}, September 2025.

\bibitem{alcock2025patent1}
G.~T.~Alcock, \emph{Systems and Methods for $\psi$-Field Drag Reduction
and Hydrodynamic Cloaking}, Patent Application, October 2025.
\end{thebibliography}

\end{document}
